\section{Signed affine horizons and smaller frame coefficients}
\label{sec:quartic23}
The starting point is the fully authenticated incoming Pell Growth 210
system described above. This round changes its cofinal horizon, and for
explicitly selected profiles also changes the frame and one orientation
error. These are semantic source changes with proofs in both directions;
they are not asserted to be pointwise bijections with the old quadratic
horizon. Later elimination, factor splitting and chart operations are exact
integer zero-set bijections.

\subsection{Recover the sign of an affine horizon}
Write $d_0=2^{19}$ and replace the horizon by
\begin{equation}\label{eq:affine-horizon}
 t=m(2z+1),\qquad m\in\{4,16\},\quad z\in\Z.
\end{equation}
Before classifying the standalone Pell equation, its divisibility row gives
$F=d_0t(2u+1)$, hence
\[
 |F|\ge d_0|t|\ge d_0m,
 \qquad F\ne0,\quad2\mid F,
 \qquad |y|\ge|F|.
\]
The conservative square frame contains the actual $F$, so
$W\ge F^2+2>258$. The retained anchor theorem therefore supplies canonical
signed main Pell pairs before recovering front signs. Since $C,F$ are even,
$D=2(C+F+1)$ is nonzero and even. Its retained power request gives
$x\ge0$ and $q=D^x$.

As $|t+1|\ge m-1$ and $m\ge4$,
\[
 |1+(t+1)y|\ge(m-1)d_0m-1>1.
\]
The terminal beta equation excludes $x=0$. Thus $x\ge1$, $|q|\ge2$.
The retained request at base $q$ and nominal index $t$ now forces $t\ge0$,
so~\eqref{eq:affine-horizon} gives $t\ge m$. This is the literal
\file{pow4.output} request at \file{index0}, not an imagined $q^{t+1}$
request. Nothing used the sign of $t$ before this step.

For $m=16$ all remaining incoming arithmetic inequalities apply unchanged.
Completeness uses the incoming cofinal horizons $t\equiv16\pmod{2^{26}}$,
which belong to the new progression $t\equiv16\pmod{32}$. This gives the
conservative 209-coordinate source used independently by the reciprocal
family.

\subsection{The smaller threshold has its own estimates}
For $m=4$, positivity gives $t\ge4$ and $4\mid t$. The constructor
inequality $E(a,a)+1\le8(a+1)^2$ is valid for every $a\ge0$, and constants
zero and one are below $t$. The same literal depth-19 trees are therefore
bounded by
\[
 L=8^{d_0-1}t^{d_0}.
\]
Pell classification at $4t+1$, the congruence $Y_n(4t+1)\equiv n\pmod{4t}$,
and the exact valuation $v_2(Y_n)=v_2(n)$ again give $n\ge d_0t$. Hence
\[
 |F|\ge(8t+1)^{d_0t-1}\ge L^t.
\]
The last comparison uses $(8t)^{d_0t-1}/L^t=8^{t-1}/t\ge1$.
These actual bounds give $|F|>2^{64}$ and $|F|>16^t$ already at $t=4$.
The remaining incoming bootstrap proves $F>0$, $y\ge F$ and
$x\ge2+(t+1)y$. Therefore $x>2^{32}$ and
\[
 2+t+x+y+F<2x.
\]
The original literal local bound
$\LL<(2x)^{2^{23}}<2^x$ is valid with these repaired hypotheses.
Also $\lcm(1,\ldots,t)\le4^{t-1}$ and $F>16^t$ supply the selector
radius. No use of an unproved $t\ge16$ remains in these steps. The radix
parity and remainder norm are unchanged.

\subsection{A cofinal binary residue construction at threshold four}
Choose arbitrarily large $t\equiv4\pmod{2^{24}}$ above the finite proof's
length and local parameters. Such a $t$ belongs to~\eqref{eq:affine-horizon},
has valuation two, and has binary bit 21 equal to zero. Put
\[
 n_0=d_0t,\quad y_0=Y_{n_0}(4t+1),\quad
 a_0=X_{n_0}(4t+1),\quad b_0=y_0/(d_0t).
\]
The incoming congruence and valuation proof gives odd integral $b_0$ and
$a_0\equiv1\pmod4$. Choose $\ell>24$ with $2^\ell>t$, and set
$r=\ell-21$. The Pell permutation lemma gives a positive odd $s$, as
large as desired, for which
\[
 (t/4)b_0Y_s(a_0)\equiv1\pmod{2^r}.
\]
Set
\[
 F=y_0Y_s(a_0),\quad X=X_{n_0s}(4t+1),\quad Z=tF,
 \quad u=(F/(d_0t)-1)/2.
\]
All growth rows hold, and $F\equiv2^{21}\pmod{2^\ell}$. Since
$t<2^\ell$ and bit 21 of $t$ is zero, $F+t$ has no binary carries;
$B=\binom{F+t}{t}$ is odd. The prime-support beta spacing, finite CRT,
even cofinal $C$, collapsed radix and final main-Pell completion now follow
in the same acyclic order as the incoming construction. Thus threshold four
has full completeness, not just a sign or growth argument.

\subsection{An optional compressed affine frame}
This further source change is used only where the endpoint certificate says
so. A literal 34-atom ledger first checks the exact incoming frame
$W=2+\sum_ja_j^2$, containing every actual base, index and proposed output
of the fifteen requests. After substituting~\eqref{eq:affine-horizon},
each $a_j$ is an integer affine form in ordinary front coordinates, not
private Pell coordinates.

For an affine form $L=c+\sum_ia_iu_i$, the integer inequality
$|a_iu_i|\le u_i^2+\lfloor a_i^2/4\rfloor$ yields
\[
 |L|\le\sum_{a_i\ne0}u_i^2+|c|+\sum_i\lfloor a_i^2/4\rfloor.
\]
The largest constant in the complete ledger is $m^2+m+1$, from $t+1$.
Replace the frame by
\begin{equation}\label{eq:compressed-frame}
 W=m^2+m+2+\sum_{\text{distinct underlying front coordinates }u_i}u_i^2.
\end{equation}
All actual affine forms stay in their power and index requests. The ledger
proves $|a_j|<W$ and the stronger bound
$|a_j|\le(3m+1)\sqrt W$, since $3m+1$ is the maximum absolute coefficient
sum including the constant, attained at $t+1$ or $F+t$. The underlying
coordinate $F$ is retained, so already $W>(d_0m)^2\ge2^{42}$.

The old squared-atom power theorem is not reused silently. Set
$s_0=\lceil\sqrt W\rceil$, $b_0=(3m+2)s_0$. Then every actual atom has
absolute value at most $b_0$, and
\[
 b_0<W,\qquad4b_0+1\le W,\qquad b_0^2\le10000W.
\]
These follow from $s_0\le2\sqrt W$, $m\le16$ and $W\ge2^{42}$;
for example $4b_0+1\le400\sqrt W+1<W$.
The exact negative-index clearing used in Theorem~\ref{thm:pooled-exactness}
now gives, for at most fifteen requests,
\[
 0\le E_*\le60b_0^{4b_0}<W^{4b_0+1}\le W^W.
\]
The modulus still satisfies
$|T_n|\ge2|A|-b_0^2-1>W^W$, because $|A|>W^W$ and
$W^W>10000W+1$. Thus $E_*=0$, giving exact powers and excluding negative
indices at bases of absolute value at least two. The same argument with one summand proves each
singleton request. No fourfold enlargement of the center is required.

\subsection{Use a difference error in orientation}
The anchor comparison first identifies the classified index $|C|$, using
$|C|<W$; the cap bounds every other classified index by $|C|$.
Replace only the orientation square
\[
 (Y_{F+t}-(F+t))^2
 \quad\hbox{by}\quad (Y_{F+t}-Y_t-F)^2,
\]
retaining $(Y_t-t)^2$ and every abscissa error. The common zero set of the
nonnegative error forms is unchanged: the retained $t$ error vanishes first,
then the difference error gives the original $F+t$ equality.

Modulo $A-1$, the changed error has absolute value below $3W$, each other
ordinate error below $2W$, and each abscissa error at most two. For thirteen
principal pairs the nonnegative error sum is therefore
\[
 <57W^2+52<4W^3<|A-1|.
\]
Its divisibility forces every error to vanish, recovering the canonical
orientation. The exceptional anchor branch is handled by the retained cap
and nonzero anchor error before invoking a comparison at an exceptional
modulus. The exact quadratic delta is propagated through the literal emitted
basis and independently checked; it is not merely an informal replacement
of a proof expression.

Completeness fixes all arithmetic-front values first, then the new $W$ and
a sufficiently large positive shared center with its private anchor witnesses.
Canonical errors are divisible by $A-1$, and the exact power identities
supply their quotients. No earlier CRT modulus depends on the eventual $W$.
Thus both the frame and orientation changes preserve the ordinary source
calibration.

\subsection{Exact elimination and coefficient splitting}
In the modified source, substitute $t=m+2mz$ everywhere and delete its
monic horizon row and coordinate. The inverse inserts that same affine
expression. This exact integral bijection includes the frame, cap,
orientation, power and growth rows, not just the visible arithmetic front.

For nonnegative integers $a_1,\ldots,a_r$ summing to 19, replace the growth
divisibility row by successive monic definitions
\[
 w_1=2^{a_1}t,\quad w_2=2^{a_2}w_1,\quad\ldots,\quad
 F=2^{a_r}w_{r-1}(2u+1).
\]
The wires have unique integer values. Forgetting and restoring them in order
are inverse zero-set maps, and every residual stays quadratic. The chosen
splits are finite certificate data, not new equation parameters. Subsequent
integer shears, translations and row additions have their recorded integral
inverses. These operations can lower coefficient height at a variable cost;
there is no general monotonicity claim.

\subsection{The frozen charts and their finite maps}
The selected endpoints use threshold four and the independently checked
compressed frame. The factor exponents and final row operations are:
\begin{center}\small
\begin{tabular}{@{}lcrr@{}}\toprule
Endpoint & Split of 19 & Residual rows & Final row additions\\\midrule
dimension209 & no extra wire &149&0\\
principal210 &$(8,11)$&150&0\\
balanced211 &$(4,8,7)$&151&4\\
low212 &$(3,5,6,5)$&152&426\\\bottomrule
\end{tabular}
\end{center}
Every metric in Section~\ref{sec:results} is obtained from full independent
coefficient collection of these actual arrays. The final balanced height
is 65537; earlier exploratory values and splits are not retained in this
table. The exact conservative $m=16$ source remains a separate 209-variable,
149-row object for reciprocal compilation, with its own complete printed
array in Appendix~\ref{sec:conservative-residuals}.


The unbounded source proofs and exact finite maps compose with the named
C222/ZF/ZFC calibration to prove~\eqref{eq:calibration} for each final
quartic. This does not assume consistency, restrict codes to externally
standard ones, or claim a fresh end-to-end formal-kernel derivation.
